% =============================================================================
% TP1 – MLOps and multivariate server-load forecasting
% Report skeleton. Max 5 pages incl. figures and tables.
% Red \todo{...} = still to fill in. Grey [Qx] tags map text to the assignment.
% Compile: pdflatex report.tex (or upload report.tex + figures/ to Overleaf).
% =============================================================================
\documentclass[10pt,a4paper]{article}

\usepackage[utf8]{inputenc}
\usepackage[T1]{fontenc}
\usepackage{lmodern}
\usepackage[margin=2cm]{geometry}
\usepackage{amsmath,amssymb}
\usepackage{graphicx}
\usepackage{booktabs}
\usepackage{tabularx}
\usepackage{xcolor}
\usepackage[font=small,labelfont=bf]{caption}
\usepackage{enumitem}
\usepackage[hidelinks]{hyperref}

\graphicspath{{../figures/}{figures/}}
\setlist{nosep,leftmargin=*}
\setlength{\parskip}{3pt}

% --- helpers -----------------------------------------------------------------
\newcommand{\todo}[1]{\textcolor{red}{\textbf{TODO:} #1}}
\newcommand{\q}[1]{\textcolor{gray}{\small[#1]}}
\newcommand{\code}[1]{\texttt{#1}}

% \fig{file}{caption}{label}: shows the image if it exists, otherwise a placeholder box.
\newcommand{\fig}[3]{%
  \begin{figure}[ht]\centering
    \IfFileExists{../figures/#1}{\includegraphics[width=0.85\linewidth]{#1}}{%
      \IfFileExists{figures/#1}{\includegraphics[width=0.85\linewidth]{#1}}{%
        \fbox{\parbox[c][3.2cm][c]{0.8\linewidth}{\centering\textcolor{red}{%
          Missing figure: \code{figures/\detokenize{#1}}}}}}}
    \caption{#2}\label{#3}
  \end{figure}}

% =============================================================================
\title{\textbf{MLOps and Multivariate Server-Load Forecasting}\\[2pt]
  \large Artificial Intelligence -- Master, TP1}
\author{Carolina \todo{surname} \quad\&\quad Jakub Weso\l{}owski}
\date{\today}

\begin{document}
\maketitle
\vspace{-1.5em}

\begin{abstract}
\noindent We forecast the one-minute-ahead CPU utilisation of a production machine from the
Alibaba Cluster Trace v2018 using its multivariate operational measurements, and study how the
model behaves when the data distribution changes after deployment. We compare a persistence
baseline, a Ridge (ARX) model and gradient-boosted trees under a strictly chronological
validation protocol, and propose a drift indicator with an alert threshold and a deployment policy.
\todo{one sentence with the main result, e.g.\ best model and test MAE.}
\end{abstract}

% =============================================================================
\section{Problem formulation}\label{sec:problem}
\q{Q1}
Let $x_t \in \mathbb{R}^p$ collect CPU, memory, memory bandwidth, cache misses, network and disk
measurements of one machine at time $t$. We learn a one-step-ahead forecaster
\begin{equation}
  \hat y_{t+1} = f\!\left(x_t, x_{t-1}, \dots, x_{t-L+1}\right), \qquad y_t = \code{cpu\_util\_percent}_t \in [0,100].
\end{equation}
\begin{itemize}
  \item \textbf{Target:} CPU utilisation of machine \code{m\_1932}.
  \item \textbf{Horizon:} one step, i.e.\ 1 minute ahead.
  \item \textbf{Sampling:} raw measurements arrive irregularly, approximately every
        \todo{verify: $\approx$10\,s}; they are averaged on a regular 1-minute grid.
\end{itemize}
\textbf{Why random shuffling is inappropriate.} Consecutive samples are strongly autocorrelated.
A shuffled split places observations from the future in the training set, so the model is
evaluated on points almost identical to ones it has already seen. The error estimate becomes
optimistically biased (temporal leakage). All splits are therefore chronological.

% =============================================================================
\section{Data and exploratory analysis}\label{sec:data}
\textbf{Source.} Alibaba Cluster Trace v2018, file \code{machine\_usage.csv}
(about 4000 machines over 8 days)~\cite{alibaba}. Memory, bandwidth, network and disk values are
normalised by Alibaba to $[0,100]$.
We extract a single machine (\code{m\_1932}) with a streaming script, so the 1.7\,GB archive never
has to be unpacked.

\textbf{Cleaning} (implemented in code, identical for every run): values outside $[0,100]$
(the documented invalid codes $-1$ and $101$) are set to missing; rows are sorted by time and
duplicate timestamps removed; samples are averaged per minute. Gaps are \emph{not} filled at this
stage so that they remain visible in the analysis.

\textbf{Descriptive statistics.} \q{Q1}
Table~\ref{tab:stats} summarises the variables after cleaning and resampling.

\begin{table}[ht]\centering\small
\caption{Descriptive statistics after cleaning (1-min resolution).}\label{tab:stats}
\begin{tabular}{lrrrrrrrr}
\toprule
Variable & count & mean & std & min & 25\% & 50\% & 75\% & max \\
\midrule
\code{cpu}     & & & & & & & & \\
\code{mem}     & & & & & & & & \\
\code{mem\_gps}& & & & & & & & \\
\code{mkpi}    & & & & & & & & \\
\code{net\_in} & & & & & & & & \\
\code{net\_out}& & & & & & & & \\
\code{disk}    & & & & & & & & \\
\bottomrule
\end{tabular}
\end{table}
\todo{fill Table~\ref{tab:stats} from \code{df.describe().T}; report \% missing per variable.}

\textbf{Temporal behaviour.} \q{Q3}
Figure~\ref{fig:timeseries} shows all variables over the 8 days.
\todo{periodicity (daily cycle?), saturation (CPU near 100\,\%?), missing intervals,
abnormal regimes.}

\fig{q3_timeseries.png}{Multivariate operational time series of machine \code{m\_1932}.}{fig:timeseries}

\subsection*{Explanatory variables}
\q{Q2}
We distinguish \emph{causal} drivers (the variable produces CPU work) from \emph{predictive}
signals (the variable co-moves with CPU without causing it). Both are legitimate inputs, because
every input is taken at time $t$ or earlier to predict $t+1$.

\begin{table}[ht]\centering\small
\caption{Explanatory variables and their expected relation to CPU load.}\label{tab:vars}
\begin{tabularx}{\linewidth}{l c l X}
\toprule
Variable & Sign & Type & Rationale \\
\midrule
\code{net\_in}   & $+$   & causal (leading)    & incoming requests reach the machine before the CPU processes them \\
\code{net\_out}  & $+$   & predictive          & outgoing traffic is a consequence of completed work \\
\code{mem}       & $+$   & mostly predictive   & co-located jobs use memory and CPU; memory stays allocated after compute ends (weak, slow) \\
\code{mem\_gps}  & $+$   & predictive          & memory bandwidth rises because the CPU is actively streaming data \\
\code{mkpi}      & $+$   & causal (efficiency) & cache misses stall the processor; stalled cycles still count as utilisation \\
\code{disk}      & $\pm$ & causal, ambiguous   & I/O-heavy batch jobs also compute ($+$), but I/O waits can leave the CPU idle ($-$) \\
lagged \code{cpu}& $+$   & autoregressive      & workloads persist over minutes; expected strongest predictor \\
\bottomrule
\end{tabularx}
\end{table}
\todo{correlation of each variable at $t$ with \code{cpu} at $t+1$; one sentence confirming or
contradicting Table~\ref{tab:vars}. Correlation $\neq$ causation.}

% =============================================================================
\section{Features and validation protocol}\label{sec:protocol}
\textbf{Lagged features.} \q{Q4}
For each variable we add lags $k \in \{1, 5, 15\}$ minutes ($x_{t-k}$, past values only). The
target is $y_{t+1}$. Increasing the lag window trades bias for variance: too few lags hide
longer dynamics (underfitting); many lags add correlated, noisy inputs and parameters and
remove usable rows (overfitting). \todo{optional: validation MAE vs.\ lag window.}

\textbf{Chronological split.} \q{Q6}
\begin{center}\small
\begin{tabular}{lll}
\toprule
Set & Time range & Use \\
\midrule
Train      & \todo{e.g.\ day 0 -- day 5} & fit model parameters \\
Validation & \todo{day 5 -- day 6}      & choose hyperparameters ($\alpha$, tree settings) \\
Test       & \todo{day 6 -- day 8}      & evaluated \emph{once}, at the end \\
\bottomrule
\end{tabular}
\end{center}
Hyperparameters are selected on the validation set only; the test set is never used for any
decision. Every learned transformation (median imputation, standardisation) is part of a
scikit-learn \code{Pipeline}~\cite{sklearn} and is therefore fitted on training data only.
\todo{number of rows dropped due to undefined lags.}

% =============================================================================
\section{Models}\label{sec:models}
\textbf{Persistence baseline.} \q{Q5}
$\hat y_{t+1} = y_t$. Any useful model must beat it.

\textbf{Ridge (ARX).} \q{Q5}
A linear model on standardised lagged inputs $z_t$, minimising squared error with an $L_2$ penalty:
\begin{equation}
  \min_{\beta_0,\beta}\ \sum_{t}\bigl(y_{t+1} - \beta_0 - z_t^{\top}\beta\bigr)^2 + \alpha\,\lVert\beta\rVert_2^2,
  \qquad
  \hat\beta = \bigl(Z^{\top}Z + \alpha I\bigr)^{-1} Z^{\top} y \;\text{(centred data)}.
\end{equation}
The penalty shrinks the weights, reduces variance and stabilises the strongly correlated lag
features; $\alpha$ is chosen on validation (\todo{selected value}).

\textbf{Gradient-boosted trees.} \q{Q7}
\code{HistGradientBoostingRegressor}~\cite{sklearn}: an additive ensemble of shallow trees, each fitted
to the residuals of the previous ones. CPU utilisation is bounded at 100\,\%; near saturation its
response to more load is flat rather than linear. Trees split on thresholds and can therefore model
plateaus, interactions and regime switches that a single linear combination cannot.
\todo{tuned hyperparameters.}

\textbf{Iterative training.} \q{Q8}
\todo{which model (boosting per iteration or small MLP); Figure~\ref{fig:learning}; diagnosis of
under/overfitting.}

\fig{q8_learning_curves.png}{Training and validation loss during training.}{fig:learning}

% =============================================================================
\section{Results}\label{sec:results}
\q{Q9}
We report MAE, RMSE, $R^2$ and the normalised error
$\mathrm{NRMSE} = \mathrm{RMSE}/\sigma_y$, where $\sigma_y$ is the standard deviation of the true
values ($\mathrm{NRMSE} < 1$ means the model beats predicting the mean).

\begin{table}[ht]\centering\small
\caption{Model comparison on validation and test data.}\label{tab:results}
\begin{tabular}{l cccc ccc}
\toprule
 & \multicolumn{4}{c}{Validation} & \multicolumn{3}{c}{Test} \\
\cmidrule(lr){2-5}\cmidrule(lr){6-8}
Model & MAE & RMSE & $R^2$ & NRMSE & MAE & RMSE & $R^2$ \\
\midrule
Persistence            & & & & & & & \\
Ridge ($\alpha=$ \,)   & & & & & & & \\
HistGradientBoosting   & & & & & & & \\
\bottomrule
\end{tabular}
\end{table}
\todo{does the ranking change with the metric? (RMSE penalises large errors more than MAE.)}

\fig{q10_pred_vs_true_time.png}{Predicted and true CPU on the test period.}{fig:predtime}

\textbf{Predicted vs.\ true.} \q{Q10}
\todo{Figure~\ref{fig:scatter}: calibration around the identity line, dispersion, outliers,
behaviour near 100\,\%.}

\fig{q10_scatter.png}{Predicted versus true CPU on the test set (identity line dashed).}{fig:scatter}

% =============================================================================
\section{Error analysis}\label{sec:errors}
\todo{residual distribution or residuals vs.\ time / vs.\ true value; where large errors occur
(spikes, saturation, missing inputs); link to the observations of Section~\ref{sec:data}.}

% =============================================================================
\section{Drift and monitoring}\label{sec:drift}
\textbf{Regime change.} \q{Q11}
\todo{real change found, or synthetic change injected (e.g.\ workload surge on \code{net\_in}
from day 6). Pre-change MAE: \_\_\,; post-change MAE: \_\_.}

\textbf{Drift indicator.} \q{Q12}
We monitor the rolling MAE over a window of \todo{e.g.\ 1\,h}. The alert threshold is
$\tau = \mu + 3\sigma$ of the rolling error measured on the \emph{validation} period
(\todo{value}). \todo{optional: input-distribution distance, e.g.\ Kolmogorov--Smirnov per feature.}

\fig{q12_drift_monitoring.png}{Rolling error with alert threshold; the vertical line marks the regime change.}{fig:drift}

% =============================================================================
\section{Deployment policy}\label{sec:policy}
\q{Q13}
\textbf{Monitored quantities:} (1) rolling MAE versus threshold $\tau$; (2) input drift score per
feature; (3) share of saturated CPU samples and of missing inputs.
\begin{itemize}
  \item \textbf{Recalibrate} (refit on a recent window, same model and hyperparameters) when the
        error alert fires but input distributions are stable.
  \item \textbf{Retrain} (re-tune hyperparameters, revisit features) when input drift persists.
  \item \textbf{Reject / fall back to persistence} when the model is worse than persistence over
        the monitoring window.
\end{itemize}
\todo{concrete numbers: window lengths, how long an alert must last before acting.}

% =============================================================================
\section{Reproducibility}\label{sec:repro}
Code, configuration and instructions are in the project repository
(\url{https://github.com/kubuswes2003/MLOps-and-multivariate-server-load-forecasting}).
Python 3.10 with pinned requirements. All settings (machine id, resampling, lags, split, random
seed $=42$) are in a single configuration file. The data are recreated from the official source
with one script, and the README gives the archive checksum. \todo{confirm: fresh clone $\rightarrow$
same figures.}

% =============================================================================
\begin{thebibliography}{9}\small
\bibitem{alibaba} Alibaba Group, \emph{Alibaba Cluster Trace Program, cluster-trace-v2018}.
  \url{https://github.com/alibaba/clusterdata}
\bibitem{sklearn} F.~Pedregosa et~al., \emph{Scikit-learn: Machine Learning in Python},
  Journal of Machine Learning Research 12, pp.~2825--2830, 2011.
\end{thebibliography}

\end{document}
